\documentclass[12pt,letterpaper]{article}
\usepackage[utf8]{inputenc}
\usepackage[spanish]{babel}
\usepackage{amsmath}
\usepackage{amsfonts}
\usepackage{amssymb}
\usepackage[margin=2.5cm]{geometry}

\begin{document}

\begin{center}
    \Large \textbf{Evaluación de Examen 2: Unidad II} \\
    \large Física para Ciencias de la Salud (005-1713) \\
    \vspace{0.5cm}
    \textbf{Estudiante:} Débora Echezuria \quad \textbf{C.I.:} 32.585.229
    \vspace{0.3cm} \\
    \large \textbf{Calificación Final: 6,5/10 puntos}
\end{center}

\vspace{0.5cm}

\section*{Análisis de Resultados}

\subsection*{1. Cinemática (Aceleración media)}
\textbf{Procedimiento y resultado:} Extrae los datos adecuadamente y plantea la fórmula de la aceleración. Sin embargo, al sustituir comete un error algebraico al escribir el numerador al revés ($0 - 0,6$ en lugar de $0,6 - 0$), pero fuerza el resultado de la fracción para que dé positivo ($4$). Además, omite las unidades en la operación y en su sección de datos indica erróneamente que la aceleración se mide en segundos ($a = 4 \text{ s}$). \\
\textbf{Veredicto:} \textbf{Parcialmente correcto} (Múltiples fallas de notación y análisis dimensional) (1/2 pts).

\subsection*{2. Dinámica (Leyes de Newton)}
\textbf{Procedimiento y resultado:} Identifica la fuerza y la masa. Sin embargo, comete un error crítico en el manejo de la fórmula: en lugar de utilizar el despeje correcto de la Segunda Ley de Newton ($a = \frac{F}{m}$), plantea una división invertida de masa entre fuerza ($a = \frac{m}{F}$). Realiza el cálculo aritmético incorrecto ($25 / 150 = 0,16$) y omite por completo las unidades de medición. \\
\textbf{Veredicto:} \textbf{Incorrecto} (Se otorgan puntos parciales por la extracción inicial de datos) (0,5/2 pts).

\subsection*{3. Trabajo Mecánico}
\textbf{Procedimiento y resultado:} Identifica perfectamente la fórmula del trabajo ($W = F \cdot d$). Efectúa la multiplicación de $400 \text{ N} \cdot 0,8 \text{ m}$ y obtiene el resultado numérico correcto de $320$, acompañado acertadamente de su unidad (Joules). \\
\textbf{Veredicto:} \textbf{Correcto} (2/2 pts).

\subsection*{4. Energía Potencial}
\textbf{Procedimiento y resultado:} Extrae correctamente los datos y aplica el producto de las tres variables según la fórmula de la energía potencial ($E_p = 1,2 \cdot 9,8 \cdot 1,5$). El resultado numérico de $17,64$ es matemáticamente exacto. No obstante, omite escribir la unidad final (Joules) al concluir el cálculo procedimental, aunque la haya anotado en su columna izquierda de datos. \\
\textbf{Veredicto:} \textbf{Parcialmente correcto} (Se aplica una leve penalización por omitir las magnitudes físicas en el resultado del cálculo) (1,5/2 pts).

\subsection*{5. Potencia Mecánica}
\textbf{Procedimiento y resultado:} Plantea correctamente la relación matemática dividiendo el trabajo efectuado entre el tiempo transcurrido ($P = \frac{W}{t}$). Desarrolla adecuadamente $75 / 60$ obteniendo el valor numérico correcto de $1,25$. Al igual que en la pregunta anterior, omite colocar la unidad correspondiente (Watts) al lado del número en su desarrollo principal. \\
\textbf{Veredicto:} \textbf{Parcialmente correcto} (Penalización por omisión de unidades) (1,5/2 pts).

\vspace{0.5cm}
\section*{Comentarios Generales y Sugerencias}
El examen denota un gran esfuerzo de la estudiante por mantener el orden, usando una estrategia metodológica útil separando la sección de \textit{Datos} de las operaciones.
\begin{itemize}
    \item Se requiere un repaso exhaustivo de las \textbf{fórmulas bases y los despejes} (notable en la Pregunta 2). Es fundamental recordar qué variable actúa como numerador y cuál como denominador.
    \item Se debe prestar máxima atención a las \textbf{unidades de medida}. Un número por sí solo en el desarrollo de física carece de significado si no está acompañado de su unidad del Sistema Internacional; asimismo, es un error grave denotar la aceleración en segundos (Pregunta 1).
    \item Se recomienda continuar puliendo los detalles teóricos para futuras evaluaciones, aprovechando el buen nivel de estructuración que ya posee.
\end{itemize}

\end{document}